\begin{textsample}{POS Dim 1 – to – Score 44.00 – to/to\_em\_7.txt}  \label{ex:f1_pos_257}
5.1.2 \textbf{Competências}, Habilidades específicas e Objetos de conhecimento de Língua Portuguesa As \textbf{competências} específicas de Linguagens e suas Tecnologias abarcam o componente curricular Língua Portuguesa.
Ressalta -se que, para o estudante desenvolver as \textbf{competências} específicas, o processo das aprendizagens essenciais ocorrerá durante os três anos do ensino \textbf{médio}, pois se entende que a \textbf{progressão} das habilidades específicas de Língua Portuguesa conduzirá à \textbf{competência}.
As 54 ( cinquenta e quatro ) habilidades específicas de Língua Portuguesa contemplam a \textbf{progressão} da aprendizagem essencial dos estudantes.
Para \textbf{tanto}, as habilidades específicas de Língua Portuguesa seguem a estrutura disposta : Para compreender a estrutura das habilidades específicas do componente curricular de Língua Portuguesa, \textbf{segue} o exemplo abaixo ( EM13LP07 ) : > ( 1 ) Analisar, ( 2 ) em textos de diferentes gêneros, marcas que expressam a posição do enunciador frente àquilo que é \textbf{dito} : uso de diferentes modalidades ( epistêmica, deôntica e apreciativa ) e de diferentes recursos gramaticais que \textbf{operam} como modalizadores ( \textbf{verbos} modais, tempos e modos verbais, expressões modais, adjetivos, locuções ou orações adjetivas, advérbios, locuções ou orações adverbiais, entonação etc. ), uso de estratégias de impessoalização ( uso de terceira pessoa e de \textbf{voz} passiva etc. ), ( 3 ) com \textbf{vistas} ao incremento da \textbf{compreensão} e da criticidade e ao manejo adequado desses elementos nos textos produzidos, considerando os contextos de produção. 1.
Verbo(s) que explicita(m) o(s) processo(s) cognitivos envolvido(s) na habilidade ; 2.
Complemento do(s) verbo(s), que explicita(m) o(s) objeto(s) de conhecimento mobilizado(s) na habilidade ; 3.
Modificadores do(s) verbo(s) ou do complemento do(s) verbo(s), que \textbf{explicitam} o contexto e/ou uma maior especificação da aprendizagem \textbf{esperada} ( visuais, verbais, sonoras, gestuais ).
Os objetos de conhecimento abarcam conteúdos ( fonologia, morfologia, variação linguística, pontuação, \textbf{progressão} temática etc. ), conceitos ( estilo, modalização, multissemiose ) e processos ( reconstrução das condições de produção, curadoria de informações, textualização, apreciação e réplica etc. ).
Compreende -se que o uso do termo “ objetos de conhecimento ” se dá como \textbf{inerente} ao próprio objeto principal de ensino e aprendizagem de Língua Portuguesa : a própria Língua/linguagem, ou seja, na busca de possibilidades de abarcá -la, é preciso mais que os conhecimentos já formalizados teoricamente, devendo considerar também as áreas de conhecimento da Linguística.
Dessa forma, visando à integração das práticas estruturantes das atividades de Língua Portuguesa, a seleção e a organização dos objetos de conhecimentos devem observar critérios \textbf{discursivos}, relativos à produção de sentidos com base em recursos e estratégias linguístico-discursivos, deixando de \textbf{lado} aquela lógica tradicional meramente estrutural ( fonologia, morfologia, sintaxe e semântica ).
E, além disso, ao abordar diversos aspectos da língua, devem -se considerar as situações enunciativas e as diversas condições de produção ( os processos ).
Daí a \textbf{compreensão} de que a opção pelo uso dessa expressão objetos de conhecimento se dá por necessidade de abarcar a amplitude da língua.
No que concerne ao componente curricular Língua Portuguesa, suas habilidades estão organizadas nos cinco Campos de Atuação, revelando a importância da \textbf{contextualização} do conhecimento escolar por meio dos campos : da vida pessoal, artístico-literário ; das práticas de estudo e pesquisa ; jornalístico-midiático ; e atuação na vida pública.
Ressalta -se que o Campo da Vida Pessoal é específico para o Ensino Médio, conforme destaca a BNCC : | Ensino Fundamental — Anos Iniciais | Ensino Fundamental — Anos \textbf{Finais} | Ensino Médio | |---|---|---| | Campo da Vida Cotidiana | | Campo da Vida Pessoal | | Campo Artístico-literário | Campo Artístico-literário | Campo Artístico-literário | | Campo das Práticas de Estudo e Pesquisa | Campo das Práticas de Estudo e Pesquisa | Campo das Práticas de Estudo e Pesquisa | | Campo da Vida Pública | Campo Jornalístico-midiático | Campo Jornalístico-midiático | | | Campo de Atuação na Vida Pública | Campo de Atuação na Vida Pública | BRASIL, 2018, p. 501.
Os Campos perpassam por toda a Área de Linguagens e suas Tecnologias, como mencionado no texto \textbf{introdutório}.
Na tabela que \textbf{segue} apresentam -se as especificidades para o componente de Língua Portuguesa : | Campo de Atuação Social | \textbf{Descrição} | |---|---| | Campo da Vida Pessoal | Pretende funcionar como espaço de articulações e sínteses das aprendizagens de outros campos postas a serviço dos projetos de vida dos estudantes.
As práticas de linguagem privilegiadas nesse campo relacionam -se com a ampliação do saber sobre si, tendo em \textbf{vista} as condições que cercam a vida \textbf{contemporânea} e as condições juvenis no Brasil e no mundo.
Está em questão também possibilitar vivências significativas de práticas colaborativas em situações de interação presenciais ou em ambientes \textbf{digitais}, inclusive por meio da articulação com outras áreas e campos, e com os projetos e escolhas pessoais dos jovens, exercendo protagonismo de forma contextualizada.
Ao se orientar para a construção do projeto de vida, a escola que acolhe as \textbf{juventudes} assume o compromisso com a formação integral dos estudantes, uma vez que promove seu desenvolvimento pessoal e social, por meio da consolidação e construção de conhecimentos, representações e valores que incidirão sobre seus processos de tomada de decisão ao longo da vida.
Dessa \textbf{maneira}, o projeto de vida é o que os estudantes almejam, projetam e redefinem para si ao longo de sua trajetória, uma construção que acompanha o ensino \textbf{médio} no desenvolvimento da(s) identidade(s), em contextos atravessados por uma cultura e por \textbf{demandas} sociais que se articulam, ora para promover, ora para constranger seus desejos.
Logo, é papel da escola \textbf{auxiliar} os estudantes a aprender a se reconhecer como sujeitos, considerando suas potencialidades e a relevância dos modos de participação e intervenção social na concretização de seu projeto de vida.
É, também, no ambiente escolar que os jovens podem experimentar, de forma mediada e intencional, as interações com o outro, com o mundo, e vislumbrar, na valorização da diversidade, oportunidades de crescimento para seu presente e futuro. | | Campo Artístico-Literário | Buscam -se a ampliação do contato e a \textbf{análise} mais fundamentada de manifestações culturais e artísticas em geral.
Está em jogo a continuidade da formação do leitor literário e do desenvolvimento da fruição.
A \textbf{análise} contextualizada de produções artísticas e dos textos literários, com destaque para os clássicos, intensifica -se no Ensino Médio.
Gêneros e formas diversas de produções vinculadas à apreciação de obras artísticas e produções culturais ( resenhas, vlogs, podcasts literários, culturais etc. ) ou a formas de apropriação do texto literário, de produções cinematográficas e teatrais e de outras manifestações artísticas ( remidiações, paródias, estilizações, videominutos, fanfics etc. ) continuam a ser considerados associados a habilidades técnicas e estéticas mais refinadas.
A escrita literária, por sua vez, ainda que não seja o \textbf{foco} \textbf{central} do componente de Língua Portuguesa, também se mostra rica em possibilidades expressivas na cultura juvenil.
O que está em questão nesse tipo de escrita não é informar, ensinar ou simplesmente comunicar.
O exercício literário inclui também a função de produzir certos níveis de reconhecimento, empatia e solidariedade e envolve reinventar, questionar e descobrir -se.
Sendo assim, ele é uma função importante em \textbf{termos} de elaboração da subjetividade e das \textbf{inter-relações} pessoais.
Nesse sentido, o desenvolvimento de textos construídos esteticamente — no âmbito dos mais diferentes gêneros — pode propiciar a exploração de emoções, sentimentos e ideias que não encontram lugar em outros gêneros não literários. | | Campo das Práticas de Estudo e Pesquisa | Mantém destaque para os gêneros e as habilidades envolvidas na leitura/escuta e produção de textos de diferentes áreas do conhecimento e para as habilidades e procedimentos envolvidos no estudo.
Ganham realce também as habilidades relacionadas à \textbf{análise}, síntese, reflexão, \textbf{problematização} e pesquisa : estabelecimento de recorte da questão ou \textbf{problema} ; seleção de informações ; estabelecimento das condições de coleta de dados para a realização de levantamentos ; realização de pesquisas de diferentes tipos ; tratamento de dados e informações ; e formas de uso e socialização dos resultados e \textbf{análises}.
Além de fazer uso competente da língua e das outras semioses, os estudantes devem ter uma atitude investigativa e criativa em relação a elas e compreender princípios e procedimentos \textbf{metodológicos} que orientam a produção do conhecimento sobre a língua e as linguagens e a formulação de regras.
No Ensino Médio, aprofundam -se também a \textbf{análise} e a reflexão sobre a língua, no que \textbf{diz} respeito à contraposição entre uma perspectiva prescritiva única, que \textbf{segue} os moldes da \textbf{abordagem} tradicional da gramática, e a perspectiva de \textbf{descrição} de vários usos da língua.
Ainda que continue a aprendizagem da norma-padrão, em função de situações e gêneros que a requeiram, outras variedades devem ter espaço e devem ser legitimadas. | | Campo Jornalístico-Midiático | Deve -se trabalhar a \textbf{compreensão} dos fatos e circunstâncias principais relatados ; perceber a impossibilidade de neutralidade absoluta no relato de fatos ; adotar procedimentos básicos de checagem de veracidade de informação ; identificar diferentes pontos de \textbf{vista} diante de questões polêmicas de relevância social ; avaliar \textbf{argumentos} utilizados e posicionar -se em relação a eles de forma ética ; identificar e denunciar discursos de ódio e que envolvam desrespeito aos Direitos Humanos ; e produzir textos jornalísticos variados, tendo em \textbf{vista} seus contextos de produção e características dos gêneros.
Eles também devem ter condições de analisar estratégias linguístico-discursivas utilizadas pelos textos publicitários e de refletir sobre necessidades e condições de consumo.
Enfatiza -se ainda mais a \textbf{análise} dos interesses que movem o campo jornalístico-midiático e do significado e das implicações do direito à comunicação e sua vinculação com o direito à informação e à liberdade de imprensa.
Também está em questão a \textbf{análise} da relação entre informação e opinião, com destaque para o fenômeno da pós-verdade, a consolidação do desenvolvimento de habilidades, a apropriação de mais procedimentos envolvidos nos processos de curadoria, a ampliação do contato com projetos editoriais independentes e a consciência de que uma mídia independente e plural é condição \textbf{indispensável} para a democracia.
Aprofundam -se também as \textbf{análises} das formas \textbf{contemporâneas} de publicidade em contexto \textbf{digital}, a dinâmica dos influenciadores \textbf{digitais} e as estratégias de engajamento utilizadas pelas empresas, sendo que as práticas que têm lugar nas redes sociais têm tratamento ampliado.
Além dos gêneros propostos para o Ensino Fundamental, são privilegiados gêneros mais \textbf{complexos} relacionados com a apuração e o relato de fatos e situações ( reportagem multimidiática, documentário etc. ) e com a opinião ( \textbf{crítica} da mídia, ensaio, vlog de opinião etc. ).
Textos, \textbf{vídeos} e podcasts diversos de apreciação de produções culturais também são propostos, a exemplo do que acontece no Ensino Fundamental, mas com \textbf{análises} mais consistentes, tendo em \textbf{vista} a intensificação da \textbf{análise} \textbf{crítica} do funcionamento das diferentes semioses. | | Campo de Atuação na Vida Pública | Está a ampliação da participação em diferentes \textbf{instâncias} da vida pública, a defesa de direitos, o domínio básico de textos legais e a discussão e o \textbf{debate} de ideias, com propostas e projetos de produção dos textos legais, sócio e historicamente situados e, em última \textbf{instância}, \textbf{baseados} nas experiências humanas, formulados com \textbf{vistas} à paz social.
Ainda no domínio das ênfases, indica -se um conjunto de habilidades que se relacionam com a \textbf{análise}, discussão, elaboração e desenvolvimento de propostas de ação e de projetos culturais e de intervenção social. | Portanto, os campos de atuação apontam a organização das práticas de linguagem ( leitura de textos, produção de textos, oralidade e \textbf{análise} linguística/semiótica ) ; são esferas \textbf{organizadoras} de modo a garantir a \textbf{contextualização} do conhecimento e articulação das habilidades, com \textbf{vistas} que essas práticas sejam situadas/empregadas em contextos significativos a partir da realidade local, regional e global.

% matched lemmas: abordagem, análise, argumento, auxiliar, basear, central, competência, complexo, compreensão, contemporâneo, contextualização, crítico, debate, demanda, descrição, digital, discursivo, dizer, esperar, explicitar, final, foco, indispensável, inerente, instância, inter-relação, introdutório, juventude, lado, maneira, metodológico, médio, operar, organizador, problema, problematização, progressão, segar, tanto, termos, verbo, vista, voz, vídeo
\end{textsample}
